\documentclass[11pt,a4paper]{article}

% ===== PACKAGES =====
\usepackage[utf8]{inputenc}
\usepackage[T1]{fontenc}
\usepackage[margin=1in]{geometry}
\usepackage{hyperref}
\usepackage{listings}
\usepackage{xcolor}
\usepackage{booktabs}
\usepackage{array}
\usepackage{longtable}
\usepackage{fancyhdr}
\usepackage{titlesec}
\usepackage{tcolorbox}
\usepackage{fontawesome5}
\usepackage{amssymb}
\usepackage{amsmath}
\usepackage{enumitem}
\usepackage{graphicx}

% ===== COLORS =====
\definecolor{codegreen}{rgb}{0.25,0.49,0.48}
\definecolor{codegray}{rgb}{0.5,0.5,0.5}
\definecolor{codepurple}{rgb}{0.58,0,0.82}
\definecolor{codeblue}{rgb}{0.13,0.29,0.53}
\definecolor{codeorange}{rgb}{0.8,0.4,0.0}
\definecolor{backcolour}{rgb}{0.95,0.95,0.95}
\definecolor{warningbg}{rgb}{1.0,0.95,0.85}
\definecolor{warningborder}{rgb}{0.9,0.6,0.2}
\definecolor{tipbg}{rgb}{0.9,0.97,0.9}
\definecolor{tipborder}{rgb}{0.3,0.7,0.3}
\definecolor{notebg}{rgb}{0.9,0.93,1.0}
\definecolor{noteborder}{rgb}{0.3,0.4,0.7}

% ===== IEC 61131-3 CODE LISTING STYLE =====
\lstdefinelanguage{IEC61131}{
    keywords={VAR, END_VAR, VAR_INPUT, VAR_OUTPUT, VAR_IN_OUT, VAR_GLOBAL, 
              TYPE, END_TYPE, STRUCT, END_STRUCT, FUNCTION, END_FUNCTION,
              FUNCTION_BLOCK, END_FUNCTION_BLOCK, PROGRAM, END_PROGRAM,
              IF, THEN, ELSE, ELSIF, END_IF, CASE, OF, END_CASE,
              FOR, TO, BY, DO, END_FOR, WHILE, END_WHILE, REPEAT, UNTIL, END_REPEAT,
              TRUE, FALSE, AND, OR, NOT, XOR, MOD, RETURN,
              ARRAY, OF, POINTER, REFERENCE, TO, AT, EXTENDS, IMPLEMENTS,
              METHOD, END_METHOD, PROPERTY, END_PROPERTY,
              INTERFACE, END_INTERFACE,
              PUBLIC, PRIVATE, PROTECTED, INTERNAL, FINAL, ABSTRACT,
              THIS, SUPER, VAR_STAT, CONSTANT, VAR_INST},
    keywordstyle=\color{codeblue}\bfseries,
    keywords=[2]{BOOL, BYTE, WORD, DWORD, LWORD, SINT, USINT, INT, UINT, 
                 DINT, UDINT, LINT, ULINT, REAL, LREAL, STRING, WSTRING,
                 TIME, LTIME, DATE, TIME_OF_DAY, TOD, DATE_AND_TIME, DT,
                 TON, TOF, TP, CTU, CTD, CTUD, R_TRIG, F_TRIG,
                 FB_Motor, FB_Valve, FB_PID, FB_Timer},
    keywordstyle=[2]\color{codepurple}\bfseries,
    keywords=[3]{ADR, SIZEOF, REF, ABS, SQRT, SIN, COS, TAN, LN, LOG, EXP,
                 SEL, MUX, LIMIT, MAX, MIN, MOVE},
    keywordstyle=[3]\color{codeorange}\bfseries,
    sensitive=false,
    morecomment=[l]{//},
    morecomment=[s]{(*}{*)},
    commentstyle=\color{codegreen}\itshape,
    stringstyle=\color{codeorange},
    morestring=[b]',
    morestring=[b]",
}

\lstdefinestyle{iecstyle}{
    language=IEC61131,
    backgroundcolor=\color{backcolour},
    basicstyle=\ttfamily\small,
    breakatwhitespace=false,
    breaklines=true,
    captionpos=b,
    keepspaces=true,
    numbers=left,
    numbersep=8pt,
    numberstyle=\tiny\color{codegray},
    showspaces=false,
    showstringspaces=false,
    showtabs=false,
    tabsize=4,
    frame=single,
    framerule=0.5pt,
    rulecolor=\color{codegray},
    xleftmargin=15pt,
    framexleftmargin=15pt,
}

\lstset{style=iecstyle}

% ===== TCOLORBOX STYLES =====
\tcbuselibrary{skins,breakable}

\newtcolorbox{warningbox}{
    colback=warningbg,
    colframe=warningborder,
    boxrule=1pt,
    left=10pt,
    right=10pt,
    top=8pt,
    bottom=8pt,
    breakable,
    before upper={\faExclamationTriangle\ \textbf{Warning:} }
}

\newtcolorbox{tipbox}{
    colback=tipbg,
    colframe=tipborder,
    boxrule=1pt,
    left=10pt,
    right=10pt,
    top=8pt,
    bottom=8pt,
    breakable,
    before upper={\faLightbulb\ \textbf{Tip:} }
}

\newtcolorbox{notebox}{
    colback=notebg,
    colframe=noteborder,
    boxrule=1pt,
    left=10pt,
    right=10pt,
    top=8pt,
    bottom=8pt,
    breakable,
    before upper={\faInfoCircle\ \textbf{Note:} }
}

% ===== HEADER/FOOTER =====
\pagestyle{fancy}
\fancyhf{}
\fancyhead[L]{\leftmark}
\fancyhead[R]{TwinCAT 3 Lecture Notes}
\fancyfoot[C]{\thepage}
\renewcommand{\headrulewidth}{0.4pt}
\renewcommand{\footrulewidth}{0.4pt}

% ===== HYPERREF SETUP =====
\hypersetup{
    colorlinks=true,
    linkcolor=codeblue,
    filecolor=magenta,
    urlcolor=cyan,
    pdftitle={Introduction to PLC Programming with TwinCAT 3},
    pdfauthor={},
    bookmarks=true,
}

% ===== TITLE FORMATTING =====
\titleformat{\section}
    {\normalfont\Large\bfseries}{\thesection}{1em}{}[{\titlerule[0.8pt]}]
\titleformat{\subsection}
    {\normalfont\large\bfseries}{\thesubsection}{1em}{}
\titleformat{\subsubsection}
    {\normalfont\normalsize\bfseries}{\thesubsubsection}{1em}{}

% ===== DOCUMENT INFO =====
\title{
    \vspace{-1cm}
    \textbf{Lecture Notes: TwinCAT 3}\\
    \Large Introduction to PLC Programming with TwinCAT 3
}
\author{}
\date{\today}

% ===== BEGIN DOCUMENT =====
\begin{document}

\maketitle

\tableofcontents
\newpage

% ============================================================
\section{Overview and Objectives}
% ============================================================

This tutorial series provides a comprehensive introduction to \textbf{TwinCAT 3 software development}, bridging the gap between traditional electrical engineering and modern software engineering practices. It is designed for both beginners and those with a software background who are new to industrial automation.

\subsection{Core Philosophy}

\begin{itemize}[leftmargin=*]
    \item \textbf{Software Engineering vs. Electrical Engineering:} Traditionally, PLC tools were rooted in electrical engineering (replacing physical relays). TwinCAT 3 represents a shift toward ``real'' software engineering.
    
    \item \textbf{Focus on Structured Text (ST):} While the IEC 61131-3 standard defines five languages (SFC, LD, FBD, IL, and ST), this series focuses exclusively on \textbf{Structured Text} because it is the most similar to high-level programming languages like C and C++.
\end{itemize}

\begin{table}[h]
\centering
\begin{tabular}{@{}llll@{}}
\toprule
\textbf{Language} & \textbf{Abbreviation} & \textbf{Type} & \textbf{Origin} \\
\midrule
Structured Text & ST & Textual & Pascal/C-like \\
Ladder Diagram & LD & Graphical & Relay logic \\
Function Block Diagram & FBD & Graphical & Signal flow \\
Instruction List & IL & Textual & Assembly-like \\
Sequential Function Chart & SFC & Graphical & State machines \\
\bottomrule
\end{tabular}
\caption{IEC 61131-3 Programming Languages}
\end{table}

\begin{notebox}
Structured Text (ST) is recommended for programmers with a software background because it supports familiar constructs like loops, conditionals, and object-oriented programming concepts that are common in languages like C, C++, and Python.
\end{notebox}

% ============================================================
\section{The TwinCAT 3 Ecosystem}
% ============================================================

The tutorial is divided into two main sections: \textbf{Basic Core Functionality} and \textbf{Advanced Ecosystem Understanding}.

\subsection{Basic Group (Parts 2--12)}

\begin{itemize}[leftmargin=*]
    \item \textbf{Environment (XAE):} Setting up the development environment and defining real-time tasks.
    
    \item \textbf{Real-Time Execution:} Investigating how a PLC defines a task and executes a program.
    
    \item \textbf{Data Handling:} Exploring standard data types and \textbf{Structures} for grouping logical variables.
    
    \item \textbf{Code Organization:}
    \begin{itemize}
        \item \textbf{Functions:} Getting data in and out of modular logic.
        \item \textbf{Function Blocks:} These are the building blocks for \textbf{Object-Oriented Programming (OOP)} in the PLC world, acting as the equivalent of classes.
    \end{itemize}
    
    \item \textbf{Standard Libraries:} Utilizing the \texttt{Tc2\_Standard} library for essential functions like timers and triggers.
\end{itemize}

\subsubsection*{Example: Basic Structured Text Variable Declaration}

\begin{lstlisting}
// PROGRAM: MAIN
// Demonstrates basic variable declaration in Structured Text
// Variables are declared within VAR...END_VAR blocks

PROGRAM MAIN
VAR
    // =====================================================
    // BOOLEAN TYPE
    // Used for TRUE/FALSE logic, digital I/O, flags
    // =====================================================
    bMotorRunning   : BOOL := FALSE;    // Motor status flag
    bStartButton    : BOOL;              // Start pushbutton input
    bStopButton     : BOOL;              // Stop pushbutton input
    bAlarmActive    : BOOL := FALSE;     // Alarm indicator
    
    // =====================================================
    // INTEGER TYPES
    // Various sizes for different numeric ranges
    // =====================================================
    nCounter        : INT := 0;          // 16-bit signed (-32768 to 32767)
    nProductCount   : UINT := 0;         // 16-bit unsigned (0 to 65535)
    nBatchNumber    : DINT := 0;         // 32-bit signed integer
    nTotalParts     : UDINT := 0;        // 32-bit unsigned integer
    nSmallValue     : SINT := 0;         // 8-bit signed (-128 to 127)
    nByteValue      : USINT := 0;        // 8-bit unsigned (0 to 255)
    
    // =====================================================
    // FLOATING POINT TYPES
    // For decimal/fractional values
    // =====================================================
    fTemperature    : REAL := 20.5;      // 32-bit float (single precision)
    fPressure       : REAL := 101.325;   // Atmospheric pressure in kPa
    fPreciseValue   : LREAL := 3.14159265358979;  // 64-bit double precision
    
    // =====================================================
    // STRING TYPE
    // Text data with specified maximum length
    // =====================================================
    sProductName    : STRING(50) := 'Widget A';   // 50 character max
    sStatusMessage  : STRING(255);                 // 255 character max
    sOperatorName   : STRING := 'Default';         // Default 80 chars
    
    // =====================================================
    // TIME TYPES
    // For delays, timers, and time-based logic
    // =====================================================
    tCycleTime      : TIME := T#100MS;   // 100 milliseconds
    tDelayDuration  : TIME := T#5S;      // 5 seconds
    tLongDelay      : TIME := T#1H30M;   // 1 hour 30 minutes
    
    // =====================================================
    // ARRAYS
    // Collections of same-type values
    // =====================================================
    aTemperatures   : ARRAY[0..9] OF REAL;         // 10 temperature readings
    aDigitalInputs  : ARRAY[1..8] OF BOOL;         // 8 digital inputs
    aBatchData      : ARRAY[0..99] OF INT;         // 100 integer values
    
END_VAR
\end{lstlisting}

\begin{lstlisting}
// Additional declaration examples showing initialization

VAR
    // Initialized array
    aSensorLimits : ARRAY[0..2] OF REAL := [0.0, 50.0, 100.0];
    
    // Constants (cannot be changed at runtime)
    CONSTANT
        c_PI         : LREAL := 3.14159265358979;
        c_MAX_SPEED  : INT := 1500;
        c_TANK_SIZE  : REAL := 1000.0;   // Liters
    END_VAR
    
    // Bit/Byte types for hardware interfacing
    byStatusByte   : BYTE := 16#00;     // 8 bits, hex notation
    wControlWord   : WORD := 16#0000;   // 16 bits
    dwDataWord     : DWORD := 16#00000000;  // 32 bits
    
END_VAR
\end{lstlisting}

\begin{tipbox}
Always initialize variables with meaningful default values. Uninitialized variables in TwinCAT default to zero/FALSE, but explicit initialization makes code more readable and prevents unexpected behavior.
\end{tipbox}

\subsection{Advanced Group (Parts 13--17)}

\begin{itemize}[leftmargin=*]
    \item \textbf{Professional Development Practices:} Implementing \textbf{Version Control} using Git and performing \textbf{Test-Driven Development (TDD)} using the \texttt{TcUnit} framework.
    
    \item \textbf{Connectivity (ADS):} Using the \textbf{Automation Device Specification} (Beckhoff's middleware) to communicate between the PLC and external applications in C\# or Linux-based C++.
    
    \item \textbf{Automation Interface:} Programmatically automating the configuration and deployment process.
\end{itemize}

\begin{table}[h]
\centering
\begin{tabular}{@{}lll@{}}
\toprule
\textbf{Topic} & \textbf{Part} & \textbf{Key Concepts} \\
\midrule
Environment Setup & 2 & XAE, XAR, Tasks \\
Data Types & 3--4 & INT, REAL, STRING, Structures \\
Functions & 5 & Inputs, Outputs, Return values \\
Function Blocks & 6--7 & OOP, Methods, Properties \\
Libraries & 8 & Tc2\_Standard, Timers, Triggers \\
Program Flow & 9 & IF/ELSE, CASE, FOR, WHILE \\
I/O Mapping & 10--11 & Hardware linking \\
State Machines & 12 & ENUM-based control \\
Version Control & 13 & Git, .gitignore \\
Testing & 14 & TcUnit, TDD \\
ADS & 15 & C\#, C++ connectivity \\
Automation & 16--17 & ITcSysManager \\
\bottomrule
\end{tabular}
\caption{Tutorial Series Overview}
\end{table}

% ============================================================
\section{Programming Logic and Instructions}
% ============================================================

The tutorial covers standard program flow instructions necessary for creating useful PLC software.

\begin{itemize}[leftmargin=*]
    \item \textbf{Conditional Logic:} Using \texttt{IF/ELSE} and \texttt{CASE} switches to manage program branches.
    \item \textbf{Iteration:} Utilizing \texttt{FOR} and \texttt{WHILE} loops for repetitive tasks.
\end{itemize}

\subsubsection*{Example: Structured Text Control Flow}

\begin{lstlisting}
// PROGRAM: ControlFlowExamples
// Demonstrates IF/ELSE statements and CASE switches in Structured Text

PROGRAM MAIN
VAR
    // Input variables
    fTemperature    : REAL := 25.0;      // Current temperature
    nMachineState   : INT := 0;          // Current state (0-4)
    bEmergencyStop  : BOOL := FALSE;     // Emergency stop button
    bAutoMode       : BOOL := TRUE;      // Automatic mode enabled
    
    // Output variables
    bHeaterOn       : BOOL;              // Heater control
    bCoolerOn       : BOOL;              // Cooler control
    bAlarm          : BOOL;              // Alarm output
    sStatusText     : STRING(100);       // Status message
    nFanSpeed       : INT;               // Fan speed percentage
    
    // Setpoints
    fTempSetpoint   : REAL := 22.0;      // Target temperature
    fHysteresis     : REAL := 2.0;       // Deadband
END_VAR

// =====================================================
// IF...THEN...ELSE Example: Temperature Control
// =====================================================

// First, check for emergency conditions (highest priority)
IF bEmergencyStop THEN
    // Emergency stop - shut everything down
    bHeaterOn := FALSE;
    bCoolerOn := FALSE;
    bAlarm := TRUE;
    sStatusText := 'EMERGENCY STOP ACTIVE';
    
ELSIF NOT bAutoMode THEN
    // Manual mode - don't control automatically
    sStatusText := 'Manual Mode - No Auto Control';
    
ELSIF fTemperature > (fTempSetpoint + fHysteresis) THEN
    // Temperature too high - activate cooling
    bHeaterOn := FALSE;
    bCoolerOn := TRUE;
    bAlarm := FALSE;
    sStatusText := 'Cooling Active';
    
ELSIF fTemperature < (fTempSetpoint - fHysteresis) THEN
    // Temperature too low - activate heating
    bHeaterOn := TRUE;
    bCoolerOn := FALSE;
    bAlarm := FALSE;
    sStatusText := 'Heating Active';
    
ELSE
    // Temperature within acceptable range
    bHeaterOn := FALSE;
    bCoolerOn := FALSE;
    bAlarm := FALSE;
    sStatusText := 'Temperature OK';
    
END_IF


// =====================================================
// CASE Statement Example: State Machine Control
// =====================================================

CASE nMachineState OF

    0:  // IDLE State
        sStatusText := 'State 0: IDLE';
        nFanSpeed := 0;
        bHeaterOn := FALSE;
        
    1:  // STARTUP State
        sStatusText := 'State 1: STARTUP';
        nFanSpeed := 25;
        // Transition logic would go here
        
    2:  // RUNNING State
        sStatusText := 'State 2: RUNNING';
        nFanSpeed := 75;
        
    3:  // COOLDOWN State
        sStatusText := 'State 3: COOLDOWN';
        nFanSpeed := 100;
        bHeaterOn := FALSE;
        
    4:  // ERROR State
        sStatusText := 'State 4: ERROR';
        nFanSpeed := 0;
        bAlarm := TRUE;
        
ELSE
    // Unknown state - go to error
    sStatusText := 'UNKNOWN STATE!';
    nMachineState := 4;  // Force to error state
    
END_CASE
\end{lstlisting}

\begin{lstlisting}
// Loop Examples: FOR and WHILE

VAR
    i               : INT;
    nSum            : INT := 0;
    aValues         : ARRAY[0..9] OF INT := [1,2,3,4,5,6,7,8,9,10];
    bSearching      : BOOL;
    nFoundIndex     : INT := -1;
    nSearchValue    : INT := 7;
END_VAR

// =====================================================
// FOR Loop: Sum all array elements
// =====================================================
nSum := 0;
FOR i := 0 TO 9 DO
    nSum := nSum + aValues[i];
END_FOR
// Result: nSum = 55

// =====================================================
// WHILE Loop: Search for a value
// =====================================================
i := 0;
bSearching := TRUE;
nFoundIndex := -1;

WHILE bSearching AND (i <= 9) DO
    IF aValues[i] = nSearchValue THEN
        nFoundIndex := i;
        bSearching := FALSE;  // Exit loop
    END_IF
    i := i + 1;
END_WHILE
// Result: nFoundIndex = 6 (0-based index where 7 is found)
\end{lstlisting}

\begin{warningbox}
WHILE loops can cause infinite loops if the exit condition is never met. Always ensure there is a guaranteed exit path, and consider adding a maximum iteration counter as a safety mechanism in production code.
\end{warningbox}

% ============================================================
\section{Hardware Interaction (I/O)}
% ============================================================

To create a functional control system, the software must interact with the physical environment.

\begin{itemize}[leftmargin=*]
    \item \textbf{Actuation:} Controlling relays, mixers, or pneumatic systems.
    \item \textbf{Feedback:} Using sensors to provide inputs back to the PLC.
    \item \textbf{Mapping:} The process of connecting software variables to physical \textbf{Inputs and Outputs (I/O)}.
\end{itemize}

\begin{table}[h]
\centering
\begin{tabular}{@{}llll@{}}
\toprule
\textbf{I/O Type} & \textbf{Direction} & \textbf{Example Device} & \textbf{Data Type} \\
\midrule
Digital Input & PLC $\leftarrow$ Field & Pushbutton, Limit switch & BOOL \\
Digital Output & PLC $\rightarrow$ Field & Relay, Indicator lamp & BOOL \\
Analog Input & PLC $\leftarrow$ Field & Temperature sensor & INT/REAL \\
Analog Output & PLC $\rightarrow$ Field & VFD speed reference & INT/REAL \\
\bottomrule
\end{tabular}
\caption{Common I/O Types in PLC Systems}
\end{table}

\begin{notebox}
In TwinCAT, I/O mapping is typically done through the System Manager by linking PLC variables (marked with \texttt{AT \%I*} for inputs or \texttt{AT \%Q*} for outputs) to physical terminal addresses. Variables can also be linked graphically in the XAE environment.
\end{notebox}

% ============================================================
\section{Scope and Limitations}
% ============================================================

While TwinCAT is vast, some ``huge'' topics are acknowledged as being outside the primary scope of this introductory tutorial but essential for a complete understanding of the ecosystem:

\begin{itemize}[leftmargin=*]
    \item \textbf{Motion Control:} Controlling motors and axes.
    \item \textbf{HMI:} Creating front-end user interfaces.
    \item \textbf{Vision \& Measurement:} Quality inspection via cameras and process analysis.
\end{itemize}

\subsubsection*{Example: Function Block Instantiation}

\begin{lstlisting}
// FUNCTION BLOCK Definition: FB_Motor
// Function Blocks are the OOP building blocks in PLC programming
// They encapsulate data (variables) and behavior (methods)
// Each instance maintains its own state - similar to objects in C++/Java

FUNCTION_BLOCK FB_Motor
VAR_INPUT
    bStart      : BOOL;     // Command to start motor
    bStop       : BOOL;     // Command to stop motor
    nSpeedRef   : INT;      // Speed reference (0-100%)
END_VAR

VAR_OUTPUT
    bRunning    : BOOL;     // Motor is running
    bFault      : BOOL;     // Motor fault detected
    nActualSpeed: INT;      // Actual speed feedback
END_VAR

VAR
    // Internal variables (private state)
    bEnabled    : BOOL;
    tRunTime    : TIME;     // Total run time
    nStartCount : UDINT;    // Number of starts
    
    // Internal timer instance
    tonRunTimer : TON;      // Timer to track run time
END_VAR

// Function Block Implementation (logic)
// This code executes each time the FB is called

// Start/Stop logic
IF bStop THEN
    bEnabled := FALSE;
ELSIF bStart AND NOT bFault THEN
    bEnabled := TRUE;
    nStartCount := nStartCount + 1;
END_IF

// Update outputs
bRunning := bEnabled AND NOT bFault;

// Simulate speed ramping (simplified)
IF bRunning THEN
    nActualSpeed := nSpeedRef;  // In reality, would ramp
ELSE
    nActualSpeed := 0;
END_IF

// Track run time
tonRunTimer(IN := bRunning, PT := T#24H);
\end{lstlisting}

\begin{lstlisting}
// PROGRAM: MAIN
// Demonstrates declaring and calling Function Block instances
// Each instance is like an "object" with its own memory/state

PROGRAM MAIN
VAR
    // =====================================================
    // FUNCTION BLOCK INSTANCES
    // Each instance is independent - has its own variables
    // Similar to creating objects from a class
    // =====================================================
    
    // Motor controller instances
    fbMotor1    : FB_Motor;     // Conveyor motor
    fbMotor2    : FB_Motor;     // Pump motor
    fbMotor3    : FB_Motor;     // Fan motor
    
    // Standard library Function Block instances
    tonDelay    : TON;          // On-delay timer
    tofDelay    : TOF;          // Off-delay timer
    ctParts     : CTU;          // Up counter
    rtStartEdge : R_TRIG;       // Rising edge detector
    ftStopEdge  : F_TRIG;       // Falling edge detector
    
    // Control variables
    bStartCmd   : BOOL;
    bStopCmd    : BOOL;
    nSpeedSetpoint : INT := 75;
    
    // Status variables
    bAnyMotorRunning : BOOL;
    nTotalStarts     : UDINT;
END_VAR

// =====================================================
// CALLING FUNCTION BLOCKS
// Must call each instance every scan to execute its logic
// Pass inputs, then read outputs after the call
// =====================================================

// Call Motor 1 - Conveyor
fbMotor1(
    bStart      := bStartCmd,
    bStop       := bStopCmd,
    nSpeedRef   := nSpeedSetpoint
);
// After call, outputs are available:
// fbMotor1.bRunning, fbMotor1.bFault, fbMotor1.nActualSpeed

// Call Motor 2 - Pump (different speed)
fbMotor2(
    bStart      := bStartCmd AND NOT fbMotor1.bFault,  // Interlock
    bStop       := bStopCmd,
    nSpeedRef   := 50
);

// Call Motor 3 - Fan (runs when either motor runs)
fbMotor3(
    bStart      := fbMotor1.bRunning OR fbMotor2.bRunning,
    bStop       := NOT (fbMotor1.bRunning OR fbMotor2.bRunning),
    nSpeedRef   := 100
);

// =====================================================
// USING STANDARD LIBRARY FUNCTION BLOCKS
// =====================================================

// Rising edge detection - detect button press
rtStartEdge(CLK := bStartCmd);
IF rtStartEdge.Q THEN
    // Start button was just pressed (rising edge)
END_IF

// On-delay timer - delay before action
tonDelay(IN := fbMotor1.bRunning, PT := T#5S);
IF tonDelay.Q THEN
    // Motor has been running for 5 seconds
END_IF

// Up counter - count parts
ctParts(CU := rtStartEdge.Q, RESET := bStopCmd, PV := 100);
// ctParts.CV = current count
// ctParts.Q = TRUE when CV >= PV

// =====================================================
// READING FUNCTION BLOCK OUTPUTS
// =====================================================
bAnyMotorRunning := fbMotor1.bRunning OR fbMotor2.bRunning OR fbMotor3.bRunning;
nTotalStarts := fbMotor1.nStartCount + fbMotor2.nStartCount + fbMotor3.nStartCount;
\end{lstlisting}

\begin{tipbox}
Function Blocks maintain their state between PLC scan cycles. This makes them ideal for implementing stateful logic like motors (running/stopped), timers (elapsed time), and counters (current count). Always call Function Block instances every scan cycle to ensure proper operation.
\end{tipbox}

% ============================================================
\section*{Quick Reference Card}
\addcontentsline{toc}{section}{Quick Reference Card}
% ============================================================

\subsection*{Common Data Types}

\begin{table}[h]
\centering
\begin{tabular}{@{}llll@{}}
\toprule
\textbf{Type} & \textbf{Size} & \textbf{Range} & \textbf{Use Case} \\
\midrule
BOOL & 1 bit & TRUE/FALSE & Digital I/O, flags \\
INT & 16 bits & -32768 to 32767 & Counters, small numbers \\
DINT & 32 bits & $\pm$2.1 billion & Large integers \\
REAL & 32 bits & $\pm$3.4e38 & Temperatures, setpoints \\
STRING & Variable & Up to 255 chars & Messages, names \\
TIME & 32 bits & T\#0ms to T\#49d & Timers, delays \\
\bottomrule
\end{tabular}
\end{table}

\subsection*{Control Flow Syntax}

\begin{lstlisting}
// IF Statement
IF condition THEN
    // code
ELSIF other_condition THEN
    // code
ELSE
    // code
END_IF

// CASE Statement
CASE nVariable OF
    0: (* code for 0 *)
    1: (* code for 1 *)
    2..5: (* code for 2-5 *)
ELSE
    (* default *)
END_CASE

// FOR Loop
FOR i := 0 TO 10 BY 1 DO
    // code
END_FOR

// WHILE Loop
WHILE condition DO
    // code
END_WHILE
\end{lstlisting}

\subsection*{Common Standard Library FBs}

\begin{table}[h]
\centering
\begin{tabular}{@{}lll@{}}
\toprule
\textbf{Function Block} & \textbf{Purpose} & \textbf{Key Output} \\
\midrule
TON & On-delay timer & Q (done), ET (elapsed) \\
TOF & Off-delay timer & Q (done), ET (elapsed) \\
TP & Pulse timer & Q (pulse active) \\
CTU & Count up & CV (count), Q (done) \\
CTD & Count down & CV (count), Q (done) \\
R\_TRIG & Rising edge & Q (edge detected) \\
F\_TRIG & Falling edge & Q (edge detected) \\
\bottomrule
\end{tabular}
\end{table}

\vfill

\begin{center}
\textit{Last Updated: \today}
\end{center}

\end{document}
